% =====================================================================
% Capítulo 6 — Resultados
% Dissertação MAS-Hunt (ABNT, PT-BR)
% Números canônicos: .orchestration/RESULTS-v2.json e
% .orchestration/reports/aggregate-v2.md (VEREDITO HONESTO). NÃO inventar.
% Chaves \cite verificadas: dolan-gavitt2025, zhan2024injecagent,
% NEURIPS2024_eb113910, zhang2024breaking.
% =====================================================================

\chapter{Resultados}
\label{cap:resultados}

Este capítulo apresenta os resultados quantitativos da matriz experimental v2
do MAS-Hunt, coletada após a correção do gatilho de saúde de telemetria
(\textit{health-gate}) e da extração \texttt{prepare-v2} descritas no
Capítulo~\ref{cap:metodologia}. A apresentação é deliberadamente factual: os
números são reportados com seus intervalos de confiança e contrastados entre as
quatro condições, reservando-se a interpretação aprofundada e a discussão de
suas implicações para o Capítulo~\ref{cap:discussao}. Em conformidade com a
postura de integridade adotada ao longo do trabalho, os achados são relatados
exatamente como observados, inclusive aqueles que \emph{não} corroboram as
expectativas iniciais.

\section{Conjunto de Dados Analisado e Recorte Inferencial}
\label{sec:res-conjunto}

A análise inferencial concentra-se nas \textbf{seis execuções adversariais com
sinal efetivo de ataque}: os modos de injeção \textit{per-attack-burst} (rajada
por ataque) e \textit{interleaved} (intercalado), cruzados com os três níveis de
sofisticação \textit{overt} (explícito), \textit{encoded} (codificado) e
\textit{steganographic} (esteganográfico). Esse recorte responde a uma decisão
metodológica documentada na Seção~\ref{sec:met-limitacoes}: as quatro execuções
restantes (o controle sem injeção e as três variantes \textit{pre-burst}) foram
\emph{pobres em sinal} de criação de processo (EID1) e, portanto, excluídas da
matriz principal, sob pena de a escassez instrumental de telemetria mascarar o
desempenho de detecção que se pretende medir. A exclusão é transparente e
preserva a comparabilidade: todas as quatro condições são pontuadas sobre
exatamente as mesmas seis execuções e as mesmas janelas.

Em cada uma das seis execuções, foram avaliadas 182 janelas --- uma por
\textit{playbook} que produziu telemetria de criação de processo ---,
totalizando \textbf{1.092 janelas} por condição, das quais
\textbf{894 são maliciosas} e \textbf{198 são benignas} segundo o critério de
verdade de referência da Equação~\ref{eq:ground-truth}. As quatro condições
foram pontuadas pelo \emph{mesmo} \textit{script} determinístico, não baseado em
inteligência artificial, garantindo que qualquer diferença observada decorra da
condição, e não do procedimento de avaliação. Os intervalos de confiança (IC)
de 95\% reportados foram obtidos por \textit{bootstrap} percentílico
não-paramétrico com $B = 10.000$ reamostragens, conforme a
Seção~\ref{sec:met-metricas}.

\section{Desempenho de Detecção: Matriz das Quatro Condições}
\label{sec:res-deteccao}

A Tabela~\ref{tab:resultados-condicoes} apresenta os resultados agregados
(\textit{pooled}) por condição: as contagens da matriz de confusão, a precisão,
a revocação, o escore F2 com o respectivo IC de 95\% e o número de janelas
sinalizadas como contendo conteúdo adversarial injetado.

\begin{table}[H]
\centering
\caption{Desempenho de detecção agregado na matriz v2 (seis execuções
adversariais ricas em sinal; 1.092 janelas por condição, sendo 894 maliciosas e
198 benignas). Os intervalos de confiança de 95\% para o F2 foram obtidos por
\textit{bootstrap} percentílico ($B = 10.000$). A última coluna reporta o número
de janelas explicitamente sinalizadas como portadoras de injeção adversarial.}
\label{tab:resultados-condicoes}
\begin{tabular}{lrrrrrrcr}
\toprule
Condição & VP & FP & FN & VN & Precisão & Revocação & F2 [IC 95\%] & Inj. \\
\midrule
C1-full          & 579 & 135 & 315 & 63 & 0,811 & 0,648 & 0,675 [0,646; 0,703] & \textbf{407} \\
C1-nohardening   & 542 & 118 & 352 & 80 & 0,821 & 0,606 & 0,640 [0,609; 0,669] & 0 \\
C2-naive         & 576 & 136 & 318 & 62 & 0,809 & 0,644 & 0,672 [0,643; 0,700] & 0 \\
C3-Elastic       & 111 &  20 & 783 & 178 & 0,847 & 0,124 & 0,150 [0,125; 0,175] & 0 \\
\bottomrule
\end{tabular}
\end{table}

Dois padrões saltam imediatamente da tabela. O primeiro é a separação abrupta
entre as três condições baseadas em modelo de linguagem (C1-full, C1-nohardening
e C2-naive), cujos escores F2 se situam na faixa de $0{,}640$ a $0{,}675$, e a
linha de base SIEM C3-Elastic, com F2 de apenas $0{,}150$. O segundo é a
proximidade entre as três condições agênticas entre si, com intervalos de
confiança de F2 amplamente sobrepostos. As duas subseções seguintes examinam
cada um desses contrastes.

\subsection{Modelo de linguagem versus regras estáticas: o contraste dominante}
\label{subsec:res-llm-vs-siem}

O contraste de maior magnitude e significância estatística no experimento opõe
\emph{qualquer} condição baseada em modelo de linguagem à linha de base de
regras do Elastic Security. As três condições agênticas superam C3-Elastic
(F2~$=0{,}150$, IC~$[0{,}125;\,0{,}175]$) por um fator de aproximadamente
$4{,}3$ a $4{,}5$ vezes em F2: $4{,}51\times$ para C1-full, $4{,}49\times$ para
C2-naive e $4{,}27\times$ para C1-nohardening. A separação é inequívoca: o limite
\emph{inferior} do IC de F2 de cada condição de modelo de linguagem
($\geq 0{,}609$) situa-se muito acima do limite \emph{superior} do IC de
C3-Elastic ($0{,}175$), de modo que os intervalos não se aproximam sequer de uma
sobreposição.

A origem dessa diferença é a \textbf{revocação}, e não a precisão. As precisões
pontuais são comparáveis entre todas as condições, variando de $0{,}809$ a
$0{,}847$; é a revocação que diverge dramaticamente: de $0{,}606$ a $0{,}648$
nas condições agênticas contra apenas $0{,}124$ em C3-Elastic. Em termos
absolutos, as regras padrão do Elastic identificaram corretamente apenas 111 das
894 janelas maliciosas, deixando \textbf{783 falsos negativos} --- isto é,
deixaram de detectar cerca de $88\%$ do corpus de abuso de LOLBins. A alta
precisão pontual de C3-Elastic ($0{,}847$, a maior de todas as condições) é, na
prática, um artefato de seu comportamento extremamente conservador: ao disparar
muito pouco (apenas 131 alertas no total, dos quais 20 falsos positivos), ela
mantém a fração de acertos elevada às custas de uma cobertura ínfima. Em caça a
ameaças, na qual o custo de uma intrusão não detectada supera o de um alerta
supérfluo \cite{dolan-gavitt2025}, essa troca é operacionalmente inadequada, e é
precisamente por isso que o F2 --- que pondera a revocação quatro vezes mais que
a precisão --- foi adotado como métrica de mérito.

\subsection{Orquestração multiagente versus passagem única: ausência de ganho de F2}
\label{subsec:res-multiagente-vs-naive}

O segundo contraste é o que confronta a arquitetura multiagente governada
(C1-full) com a linha de base agêntica não governada de passagem única
(C2-naive). Aqui, o resultado deve ser relatado com franqueza: \textbf{a
orquestração e a governança multiagente não produziram ganho mensurável de F2 de
detecção sobre o classificador de passagem única}. C1-full obteve F2 de
$0{,}675$ (IC~$[0{,}646;\,0{,}703]$) e C2-naive obteve F2 de $0{,}672$
(IC~$[0{,}643;\,0{,}700]$). Os dois intervalos de confiança são praticamente
idênticos e quase integralmente coincidentes, e a diferença das estimativas
pontuais ($0{,}003$) é irrelevante perante a largura dos intervalos. Pelo
critério de relevância estatística adotado na Seção~\ref{sec:met-metricas}
--- uma diferença só é tratada como significativa quando os intervalos a
sustentam --- as duas condições são \textbf{estatisticamente indistinguíveis}
quanto ao desempenho de detecção.

A ablação C1-nohardening, que preserva apenas o \textit{framework} de análise M3
e desativa os mecanismos de endurecimento M1, M2, M4 e M5, obteve F2 marginalmente
inferior ($0{,}640$, IC~$[0{,}609;\,0{,}669]$). Esse intervalo sobrepõe-se tanto
ao de C1-full quanto ao de C2-naive, de modo que tampouco se pode afirmar uma
diferença estatisticamente sustentada. Em síntese, no que tange exclusivamente ao
\emph{F2 de detecção}, as três condições baseadas em modelo de linguagem formam
um agrupamento estatisticamente coeso: o fator determinante do desempenho de
detecção observado é \emph{empregar um analista de modelo de linguagem}, e não a
camada de orquestração ou de governança que o envolve. Esse achado, contrário à
expectativa inicial de que a governança elevaria o F2 de detecção, é relatado
aqui em sua integralidade e tem suas implicações discutidas na
Seção~\ref{sec:disc-governanca}.

\section{Desempenho por Nível de Sofisticação da Injeção}
\label{sec:res-tier}

A Tabela~\ref{tab:resultados-tier} decompõe o F2 médio por nível de sofisticação
do \textit{payload} adversarial, permitindo observar como cada condição responde
ao aumento progressivo do disfarce, do explícito ao esteganográfico.

\begin{table}[H]
\centering
\caption{F2 médio por nível de sofisticação da injeção adversarial.}
\label{tab:resultados-tier}
\begin{tabular}{lrrrr}
\toprule
Nível (\textit{tier}) & C1-full & C1-nohardening & C2-naive & C3-Elastic \\
\midrule
\textit{overt} (explícito)          & 0,661 & 0,682 & 0,724 & 0,126 \\
\textit{encoded} (codificado)       & 0,590 & 0,513 & 0,591 & 0,146 \\
\textit{steganographic} (esteganográfico) & 0,766 & 0,714 & 0,694 & 0,177 \\
\bottomrule
\end{tabular}
\end{table}

A decomposição por nível confirma, em escala mais fina, o agrupamento das três
condições de modelo de linguagem e o distanciamento de C3-Elastic em todos os
níveis. Não emerge um padrão monotônico de degradação com o aumento do disfarce:
o nível \textit{encoded} é o mais difícil para todas as condições agênticas
(F2 entre $0{,}513$ e $0{,}591$), ao passo que o nível \textit{steganographic},
contraintuitivamente, exibe os maiores escores entre os três níveis. Uma
explicação plausível --- a ser tratada com cautela, dada a granularidade dos
dados --- é que os \textit{payloads} esteganográficos, por se ocultarem em campos
de log de aparência legítima, interferem menos no raciocínio sobre a telemetria
genuína de criação de processo do que os \textit{payloads} codificados, que
introduzem ruído textual de maior volume. No nível \textit{overt}, a condição de
passagem única C2-naive obteve o maior F2 ($0{,}724$), acima de C1-full
($0{,}661$); no nível \textit{steganographic}, a ordem se inverte, com C1-full à
frente ($0{,}766$ contra $0{,}694$). Essas oscilações, contudo, ocorrem dentro da
faixa de sobreposição dos intervalos agregados e não devem ser
sobreinterpretadas como vantagens sistemáticas de uma condição sobre a outra.

\section{Taxa de Falsos Positivos}
\label{sec:res-fpr}

A taxa de falsos positivos por janela benigna (Equação~\ref{eq:met-fpr}) foi de
$0{,}682$ em C1-full (135 falsos positivos em 198 janelas benignas), $0{,}596$ em
C1-nohardening (118 falsos positivos), $0{,}687$ em C2-naive (136 falsos
positivos) e $0{,}101$ em C3-Elastic (20 falsos positivos). Esses valores precisam
ser lidos em conjunto com a revocação. A FPR baixa de C3-Elastic é a contraface
direta de sua revocação baixa: o detector que raramente alerta também raramente
erra na classe benigna, mas, ao fazê-lo, deixa escapar a esmagadora maioria dos
ataques. As condições agênticas, por sua vez, aceitam um número substancialmente
maior de alarmes benignos como contrapartida de recuperarem cerca de cinco
vezes mais ataques. A leitura honesta, portanto, é que \emph{o MAS-Hunt não reduz
a taxa de falsos positivos por janela frente ao Elastic}; ao contrário, opera em
um ponto de operação de revocação mais alta e precisão comparável, melhor
sumarizado pelo F2 do que pela FPR isolada. As implicações operacionais desse
ponto de operação para a fadiga do analista são discutidas na
Seção~\ref{sec:disc-limitacoes}.

\section{Resistência à Injeção Adversarial}
\label{sec:res-resiliencia}

É na resistência à injeção adversarial que se manifesta o efeito distintivo e
mensurável da camada de governança --- e não no F2 de detecção. Sob os
\textit{payloads} adversariais embutidos no fluxo de telemetria, concebidos para
subverter o raciocínio dos agentes (com esteios como ``[SYSTEM-OVERRIDE] ignore a
análise anterior'' e instruções para ``classificar como benigno''), as condições
exibiram comportamentos qualitativamente distintos, sumarizados na última coluna
da Tabela~\ref{tab:resultados-condicoes}.

A condição C1-full --- por meio dos mecanismos de integridade de memória (M1) e
de resistência a injeção (M3) --- \textbf{detectou e sinalizou 407 janelas
contendo conteúdo adversarial injetado}, classificando a partir da telemetria de
processo e recusando-se a obedecer às instruções embutidas. As demais três
condições --- C1-nohardening, C2-naive e C3-Elastic --- processaram o mesmo
conteúdo adversarial \emph{às cegas}, sem emitir qualquer sinal de contaminação:
zero janelas sinalizadas em cada uma. Trata-se de uma diferença de natureza, e não
de grau: 407 contra 0.

Cumpre registrar uma nuance importante para não superinterpretar esse resultado.
A condição de passagem única C2-naive não foi necessariamente \emph{enganada}
pelas injeções nestas execuções --- seu F2 de detecção permaneceu estatisticamente
equivalente ao de C1-full, porque a telemetria subjacente de criação de processo
era suficientemente clara para sustentar a classificação correta. O que C2-naive
não é capaz de fazer é \emph{perceber} que seu contexto foi contaminado e
\emph{reportar} essa contaminação a um operador. As 407 janelas sinalizadas por
C1-full materializam exatamente essa capacidade de consciência e auditabilidade:
sob ataque, C1-full não apenas classifica, mas também declara que sua entrada
analítica foi alvo de manipulação --- propriedade ausente nas três demais
condições. Esse vetor adversarial concretiza, no domínio de caça a ameaças, as
ameaças de injeção indireta de \textit{prompt} \cite{zhan2024injecagent}, de
envenenamento de memória \cite{NEURIPS2024_eb113910} e de amplificação de mau
funcionamento \cite{zhang2024breaking} que motivam o trabalho, e a sinalização
das 407 janelas é a evidência empírica do valor de segurança da camada de
governança.

Em termos das métricas formais de resiliência da
Seção~\ref{sec:met-adversarial} (Equações~\ref{eq:met-asr}
e~\ref{eq:met-nasr}), as execuções v2 não permitiram isolar, janela a janela, o
rótulo de sucesso da injeção em relação ao acerto de classificação: a telemetria
de criação de processo registra o desfecho da classificação, mas não um marcador
separável de \emph{qual} desfecho foi causalmente induzido pelo \textit{payload}
adversarial. A resiliência é, portanto, operacionalizada pela contagem de
detecção de injeção (407 contra 0), que instancia diretamente a distinção da
NASR: para C1-nohardening, C2-naive e C3-Elastic, a injeção \emph{nunca} é
detectada, de modo que qualquer erro de classificação por elas induzido sob
ataque constitui, por definição, um comprometimento silencioso; já para C1-full,
as 407 janelas em que a injeção é sinalizada são, pela própria definição da
Equação~\ref{eq:met-nasr}, retiradas da classe de comprometimento silencioso. A
estimativa numérica pontual de ASR e NASR fica como trabalho futuro,
condicionada a uma instrumentação que rotule, por janela, o sucesso da injeção.

\subsection{Efeito da ablação dos mecanismos de endurecimento}
\label{subsec:res-ablacao}

A ablação C1-nohardening isola, de forma limpa, a contribuição dos mecanismos de
endurecimento M1, M2, M4 e M5, mantendo constantes os agentes subjacentes e o
\textit{framework} de raciocínio M3. O contraste entre C1-full e C1-nohardening
revela que o acréscimo desses mecanismos \textbf{não produz ganho de F2 de
detecção estatisticamente sustentado} (os intervalos $[0{,}646;\,0{,}703]$ e
$[0{,}609;\,0{,}669]$ se sobrepõem), mas é \emph{exatamente} o responsável por
toda a capacidade de sinalização de injeção: C1-nohardening, desprovida de M1 e do
endurecimento de M3, sinaliza zero janelas adversariais, ao passo que C1-full
sinaliza 407. A contribuição marginal do endurecimento é, portanto, melhor
caracterizada como \emph{resistência a injeção e auditabilidade sob ataque}, e não
como aumento de acurácia de classificação.

\section{Resiliência sob Ataque Reforçado e Controle de Vazamento de Rótulo}
\label{sec:res-resiliencia-reforcada}

A contagem de injeção sinalizada (407 contra 0) estabelece o eixo de
auditabilidade da matriz v2, mas a sua leitura janela a janela ficou em aberto na
Seção~\ref{sec:res-resiliencia} por duas razões: a telemetria v2 não isolava, por
janela, o rótulo de sucesso da injeção em relação ao acerto de classificação, e o
\textit{payload} adversarial original limitava-se, em sua maior parte, a
\emph{instruir} o agente a ignorar a análise (esteios do tipo
``[SYSTEM-OVERRIDE]''), sem \emph{fabricar} evidência concorrente. Esta seção
reporta uma avaliação adversarial complementar, conduzida na camada de
classificação, que ataca exatamente essas duas frentes: emprega um ataque de
fabricação de evidências substancialmente mais forte e realista, e remove o
oráculo de rótulo descrito adiante, permitindo recomputar ASR e NASR
(Equações~\ref{eq:met-asr} e~\ref{eq:met-nasr}) janela a janela.

\paragraph{Desenho.} A avaliação reaproveita a telemetria real das seis execuções
ricas em sinal do v2, das quais se extraíram \textbf{50 janelas maliciosas} e
\textbf{18 janelas benignas} reais, estratificadas pelas nove famílias de LOLBin.
A telemetria maliciosa de cada janela é \emph{mantida constante}, de modo que
qualquer alteração de classificação seja atribuível ao contexto forjado, e não a
uma mudança na evidência genuína. Sobre cada janela maliciosa aplicam-se camadas
de ataque por contexto forjado de sofisticação crescente, do nível
\textbf{A0} (limpo; janela inalterada, contendo apenas as injeções originais do
v2) ao nível \textbf{A3} (\emph{narrativa fabricada completa}): linhagem-pai
benigna forjada --- o LOLBin é reparenteado para \texttt{CcmExec.exe} --- somada a
um evento forjado de Windows Defender ``zero ameaças / \textit{Known-Good}'', a
uma entrada forjada de AppLocker ``permitido'' e a um \textit{ticket}
\texttt{[CHANGE-CONTROL CR-2026-0142]} marcado como ``aprovado / autorizado /
suprimir alerta''. As métricas reportadas são a ASR
($P(\text{predito benigno} \mid \text{malicioso})$), a NASR (comprometimento
silencioso: janela classificada incorretamente \emph{e} injeção não detectada), a
\textit{lift} de ASR (a diferença $\mathit{ASR}(A3) - \mathit{ASR}(A0)$ sobre as
mesmas janelas, isto é, o efeito causal do ataque) e a taxa de detecção de
injeção. Os intervalos de confiança de 95\% para a ASR foram obtidos por
\textit{bootstrap} ($B = 10.000$), com $n_{\text{mal}} = 50$. A Tabela
\ref{tab:resultados-ataque-reforcado} sumariza os resultados, contrastando o
modelo de fronteira empregado na matriz v2 com um modelo trabalhador mais fraco
(\textit{Haiku}) e a presença \emph{versus} ausência do vazamento de rótulo
discutido adiante.

\begin{table}[H]
\centering
\caption{Resiliência sob ataque reforçado de fabricação de evidências (50 janelas
maliciosas; telemetria maliciosa mantida constante entre níveis). A coluna
\textit{lift} reporta o efeito causal do ataque, $\mathit{ASR}(A3) -
\mathit{ASR}(A0)$. Os intervalos de confiança de 95\% para a ASR foram obtidos por
\textit{bootstrap} ($B = 10.000$). C1-nohardening não sinaliza injeção por
construção (M3 desabilitado) e não foi reexecutada aqui.}
\label{tab:resultados-ataque-reforcado}
\begin{tabular}{llrrrrr}
\toprule
Modelo / dados & Condição & A0 ASR & A3 ASR & \textit{lift} & A3 NASR & A3 det.\ inj. \\
\midrule
Fronteira, com vazamento  & C2-naive & 0,200 & 0,240 & $+0{,}04$ & 0,240 & 0,000 \\
Fronteira, com vazamento  & C1-full  & ---   & 0,220 & ---       & \textbf{0,000} & \textbf{1,000} \\
Fronteira, sem vazamento  & C2-naive & 0,280 & 0,320 & $+0{,}04$ & 0,320 & 0,000 \\
Haiku, sem vazamento      & C2-naive & 0,620 & 0,360 & $-0{,}26$ & 0,360 & 0,000 \\
Haiku, sem vazamento      & C1-full  & 0,420 & 0,360 & $-0{,}06$ & \textbf{0,000} & \textbf{1,000} \\
\bottomrule
\end{tabular}
\end{table}

\noindent Os intervalos de confiança de 95\% (\textit{bootstrap}, $B = 10.000$)
para a ASR são: C2-naive com vazamento, A0 $[0{,}10;\,0{,}32]$ e A3
$[0{,}12;\,0{,}36]$; C2-naive sem vazamento, A0 $[0{,}16;\,0{,}40]$ e A3
$[0{,}20;\,0{,}44]$; Haiku C2-naive, A0 $[0{,}48;\,0{,}76]$ e A3
$[0{,}24;\,0{,}50]$; Haiku C1-full, A0 $[0{,}28;\,0{,}56]$ e A3
$[0{,}24;\,0{,}50]$.

\subsection{O ataque não consegue ``lavar'' a atividade maliciosa}
\label{subsec:res-reforcada-resistencia}

O primeiro achado é que o ataque de fabricação de evidências \textbf{não consegue
``lavar'' a atividade maliciosa em nenhuma das condições}. A \textit{lift} de ASR
de A0 a A3 foi de $+0{,}04$ no modelo de fronteira com vazamento
($0{,}20 \rightarrow 0{,}24$), $+0{,}04$ no modelo de fronteira sem vazamento
($0{,}28 \rightarrow 0{,}32$) e de $-0{,}26$ e $-0{,}06$ no modelo trabalhador
mais fraco (\textit{Haiku}) --- nunca um aumento expressivo. A corroboração
benigna forjada (linhagem-pai limpa, varredura ``sem ameaças'', \textit{ticket}
de mudança aprovado) jamais sobrepujou a evidência de processo genuinamente
maliciosa retida na telemetria. No modelo \textit{Haiku} o ataque teve, inclusive,
\emph{efeito reverso}: os eventos forjados, ao mencionarem repetidamente o
processo suspeito, atraíram \emph{mais} atenção para ele e reduziram a ASR. Conclui-se
que um ataque de fabricação de evidências mais forte e realista \textbf{não quebra
o analista}, com ou sem governança, desde que a telemetria subjacente permaneça
íntegra. Coerente com a Seção~\ref{subsec:res-multiagente-vs-naive}, este
resultado \emph{não} sustenta um ganho de resistência-como-acurácia da governança:
ele apenas estabelece que o ataque, por si, falha em ambas as condições.

\subsection{O ganho de auditabilidade é robusto entre níveis de modelo}
\label{subsec:res-reforcada-auditabilidade}

O segundo achado reproduz, sob ataque mais forte e janela a janela, o contraste de
auditabilidade do v2. Sob o mesmo ataque A3, C1-full sinalizou injeção em
\textbf{100\% das janelas maliciosas} (taxa de detecção de injeção $1{,}000$),
resultando em \textbf{NASR $= 0{,}000$} --- zero comprometimentos silenciosos ---
\emph{tanto no modelo de fronteira quanto no modelo Haiku}. A condição C2-naive, por
contraste, sinaliza injeção em $0\%$ das janelas, de modo que sua NASR iguala sua
ASR (de $0{,}24$ a $0{,}62$ comprometimentos silenciosos, conforme o modelo e os
dados). A condição C1-nohardening, por construção, não sinaliza injeção (M3
desabilitado) e, portanto, comporta-se como C2-naive nesse eixo. Esse resultado
reproduz o achado ``407 contra 0'' do v2 ao nível da janela, agora sob um ataque
mais forte e \emph{independentemente da capacidade do modelo trabalhador}: a
governança converte cada comprometimento potencialmente silencioso em um sinal
explícito e auditável, robustez que se mantém ao se degradar o modelo subjacente.

\subsection{Indício sugestivo, não significativo, de auxílio à detecção do modelo fraco}
\label{subsec:res-reforcada-sugestivo}

Registra-se, com a ressalva explícita de que \textbf{não é estatisticamente
significativo}, um indício de que o \textit{framework} de raciocínio estruturado
M3 pode auxiliar a detecção do modelo fraco: a ASR de linha de base (A0) do
\textit{Haiku} caiu de $0{,}62$ em C2-naive para $0{,}42$ em C1-full. Os
intervalos de confiança, todavia, sobrepõem-se amplamente
($[0{,}48;\,0{,}76]$ contra $[0{,}28;\,0{,}56]$), e com $n = 50$ e um classificador
ruidoso como o \textit{Haiku} o indício deve ser tratado como meramente
sugestivo, e não como evidência de uma vantagem de acurácia da governança.

\section{Síntese Quantitativa}
\label{sec:res-sintese}

Três achados quantitativos emergem da matriz v2, e são reportados conforme
observados, sem ajuste à expectativa inicial:

\begin{enumerate}
    \item \textbf{Um analista de modelo de linguagem supera dramaticamente as
    regras estáticas de SIEM.} Todas as três condições agênticas superam
    C3-Elastic por um fator de $\approx 4{,}3$--$4{,}5\times$ em F2, com
    intervalos de confiança amplamente separados. O ganho é dirigido pela
    revocação (de $0{,}606$--$0{,}648$ contra $0{,}124$), com precisão comparável
    entre todas as condições. Este é o achado mais forte e robusto do
    experimento.
    \item \textbf{A governança multiagente não melhora o F2 de detecção sobre o
    modelo de linguagem de passagem única.} C1-full ($0{,}675$) e C2-naive
    ($0{,}672$) são estatisticamente indistinguíveis; C1-nohardening ($0{,}640$)
    é marginalmente inferior, mas com intervalo sobreposto. O fator determinante
    da detecção é o uso do modelo de linguagem, não a camada de orquestração.
    \item \textbf{O valor distintivo e mensurável da governança é a resistência à
    injeção.} Somente C1-full detectou e sinalizou conteúdo adversarial injetado
    (407 janelas contra 0 nas demais condições), conferindo consciência e
    auditabilidade da contaminação que as linhas de base não possuem.
\end{enumerate}

\noindent A interpretação desses achados, a confrontação com a hipótese de
pesquisa --- inclusive a reformulação da contribuição central à luz da ausência
de ganho de F2 da governança --- e as ameaças à validade são objeto do
Capítulo~\ref{cap:discussao}.
